\scanpage{135}
% archive foliation: 134
\uline{Lemme} Soit $S$ un préschém. loc noeth., et pour \unsure{tout|tt} $s\in S$, soit $\deleted{\unsure{X_s}}$ $X(s)$ un préschém. quasi fini \uline{étale} \uline{séparé} sur $\operatorname{Spec}(\mathcal O_s)$, et pour $t\in S$, $s\in\overline t$, soit $\varphi_{st}:X(t)\unsure{-|\to}X(s)$ rendant cartésien le diagramme ci-contre.
\diagram{\begin{tikzcd}[column sep=large]
X(s) \arrow[d] & X(t) \arrow[l] \arrow[d] \\
\operatorname{Spec}(\mathcal O_s) & \operatorname{Spec}(\mathcal O_t) \arrow[l]
\end{tikzcd}}
On suppose que $\varphi_{ss}=\text{identité}$
\[
\varphi_{st}\varphi_{tu}=\varphi_{su}.
\]
\deleted{\unsure{des}}
\inserted{[On suppose que pour tt $s\in S$, il existe \unsure{une définition de} \deleted{\unsure{la}} $\mathcal X$ de $X(s)$ sur un voisinage ouvert $U$ de $s$, \deleted{\unsure{le groupe}} pour tt $t\in U$, un hom \deleted{$\unsure{X_t}$} $\mathcal X\times_U\operatorname{Spec}(\mathcal O_t)\to X(t)$, ces hom étant compatibles avec les \unsure{$\varphi$}.]}
\diagram{\begin{tikzcd}[column sep=small,font=\small]
X(s) \arrow[r,no head] \arrow[d] & \mathcal X \arrow[d] & \mathcal X\times_U\operatorname{Spec}(\mathcal O_t) \arrow[l] \arrow[r] \arrow[d] & X(t) \\
\operatorname{Spec}(\mathcal O_s) \arrow[r] & U & \operatorname{Spec}(\mathcal O_t) \arrow[l] &
\end{tikzcd}}
Alors il existe un préschém. étale \inserted{quasi-fini séparé} sur $S$, soit $X$, et des \inserted{$\mathcal O_s$-}isom
\[
\varphi_s:X\times_S\operatorname{Spec}(\mathcal O_s)\unsure{-|\to}X(s)
\]
rendant commutatifs les diagrammes qu'on pense. $X$ est unique à isom unique près.

Ce dernier fait résulte d'un énoncé plus général sur les catégories des préschémas de type fini sur $S$, qui donne un foncteur pleinement fidèle. \unsure{Je vais pour du moins ramener au local} et supposer $S$ affine, donc noeth. \unsure{le prouver} par récurrence noethérienne, en supposant l'énoncé prouvé pour les sous-préschémas fermés de $S\ne S$.
